\textbf{From local opportunity to integrated post-prefill decode outcomes.} Local acceptance, isolated drafting, and target-forward batching do not by themselves establish integrated gains. Arrows mark evaluation boundaries, not a guarantee of gain. The frozen greedy token-ID contract excludes the tested RTX 5070 Ti Laptop batched path; B570 serialization improves agreement without qualifying the full workload. Blocked timings remain diagnostic. For the qualified 358H NPU path, full-prefix replay and proposal share one synchronous OpenVINO call, and the isolated draft advantage does not survive this measured state interface. $D_{\mathrm{real}}$, $R_{\mathrm{target}}$, and $X$ normalize measured draft, complete target, and residual costs by matched target-only token cost; $C$ counts all committed output, including correction/bonus tokens. M4\textquotesingle s separate frozen studies on the same tested Metal stack yield opposite exact regimes. Workload and session conditions differ, and neither diagnosis is assigned to the contrast. The measured mechanisms are configuration-scoped and non-exhaustive; timing excludes loading and initial prompt prefill. The RTX interval spans measured contexts, not a CI; $V_2^{\mathrm{micro}}$ is not complete $R_{\mathrm{target}}$. The 358H throughput values are isolated medians; their 95\% repetition-bootstrap CIs are [47.10, 48.40] (NPU) and [30.34, 31.36] (CPU). The cost fraction is a ratio of arithmetic totals of speculative post-prefill time; speedup is a prompt geometric mean. The 358H NPU path is 90/90 exact, with speedup 95\% CI [0.333, 0.369]. M4 24Q and HumanEval30 fixed-$k=2$ are 72/72 and 90/90 exact; their prompt-GM speedup 95\% CIs are [0.809, 0.877] and [1.351, 1.422].
